\subsection{Actual pure classical multiplicity objects via geometric Satake}
\begin{finding}[\fid{F20}: Schur summands have pure multiplicity spaces; their cycle action need not be positive]\label{f:satake}
Let $G=\GL_r/\F_q$ be split, with $\ell\nmid q$, and use the symmetric geometric Satake category with $\overline\Q_\ell$ coefficients. For a dominant partition $\mu$ of length at most $r$, let $\IC_\mu$ be the natural unnormalized Weil intersection complex of its Schubert variety, perverse of weight
\[
 D_\mu=\dim\Gr_\mu=\langle2\rho,\mu\rangle
 =\sum_{i=1}^r(r+1-2i)\mu_i.
\]
For $\nu\vdash d$, a primitive Young idempotent $e_\nu\in\Q[S_d]$ acts, through the parity-corrected symmetric commutativity, on
$K=\IC_\mu^{\star d}$. Its image $P_{\nu,\mu}$ is a pure perverse Weil sheaf of weight $dD_\mu$. The canonical multiplicity Weil module
\[
 M_{\lambda;\nu,\mu}
 =\Hom_{\mathrm{Sat}_{\overline{\F}_q}}
 (\IC_\lambda,P_{\nu,\mu})
\]
is pure of weight $dD_\mu-D_\lambda$ and has dimension
\[
 \dim M_{\lambda;\nu,\mu}
 =[S_\lambda(\overline\Q_\ell^r):
   S_\nu(S_\mu(\overline\Q_\ell^r))].
\]
On the product Grassmannian for $\GL_m\times\GL_n$, apply the Schur idempotent for $\nu\vdash d$ to the $d$th convolution power of
$\IC_{(1),m}\boxtimes\IC_{(1),n}$. Its multiplicity module at
$\IC_{\lambda,m}\boxtimes\IC_{\mu,n}$ has dimension $g(\lambda,\mu,\nu)$ and is pure of weight
\[
 d(m+n-2)-D_{\lambda,m}-D_{\mu,n}.
\]
These are ordinary geometric realizations of the classical coefficients, including rectangular cases within the dimension bounds. They do not by themselves identify the nonstandard quantum bases.

Even in this construction the symmetric-group action on a geometric multiplicity space need not permute a basis. For $G=\GL_2$, $\mu=(1)$ and $d=2$, the determinant multiplicity line is $\overline\Q_\ell(-1)$ and the transposition acts by $-1$. Thus no basis of that line is an $S_2$-permutation basis.
\end{finding}
\begin{proof}
The literature inputs are geometric Satake, its finite-field convolution/symmetry, and ordinary weight theory. We use the published version of Richarz \cite[Theorem~3.1, p.~214; proof pp.~226--227; Proposition~4.1, p.~227; Theorem~5.1, p.~230; Lemma~5.4 and its proof, p.~232]{Richarz}. The tensor equivalence is also Mirkovi\'c--Vilonen, Theorem~7.3 \cite{MV}. These theorems are cited, not re-proved here.

The twisted product of $d$ Schubert varieties is projective and is smooth-locally the product for purposes of its intersection complex. Its natural Weil IC is normalized to restrict to $\overline\Q_\ell[dD_\mu]$ on the dense twisted orbit. With that normalization, rather than only a geometric IC isomorphism, it is the twisted external product of the $\IC_\mu$ and is pure of weight $dD_\mu$; this is the IC identification in the proof of Richarz's Lemma~5.4, together with BBD~5.3.2, second proof (which covers a nonaffine open immersion). Its proper convolution pushforward is pure of that same complex weight (BBD, Remark~5.4.9) and is perverse (Richarz, Theorem~3.1). This proves the assertion about $K$.

The symmetry constructed over $\F_q$, with the component-parity sign specified in the cited proof, gives an $S_d$ action by Weil endomorphisms. In the abelian perverse category the idempotent splits; its image is a direct summand, hence remains pure by BBD~5.3.1. Under the symmetric monoidal Satake equivalence, $\IC_\mu$ corresponds to $S_\mu(\overline\Q_\ell^r)$ and $K$ to its $d$fold tensor power. In Schur--Weyl decomposition, the primitive idempotent $e_\nu$ has rank one on the Specht factor $[\nu]$ and vanishes on the other factors. Its image therefore corresponds to one copy of the Schur functor $S_\nu(S_\mu(\overline\Q_\ell^r))$. A central isotypic idempotent would instead give $\dim[\nu]$ copies; that is why the primitive condition matters.

Here is the Frobenius argument for the multiplicity weight. Geometric semisimplicity gives the canonical evaluation isomorphism
\[
 \bigoplus_\lambda\IC_\lambda\otimes
 \Hom(\IC_\lambda,\overline P_{\nu,\mu})
 \longrightarrow\overline P_{\nu,\mu}.
\]
It is the canonical isotypic decomposition, not an arbitrarily chosen decomposition. Every support and natural $\IC_\lambda$ is defined over the split base $\F_q$, so Frobenius preserves each summand. Give its Hom factor the conjugation Frobenius induced by the two fixed Weil structures. Evaluation is then Weil-equivariant, and the isotypic summand is a pure perverse Weil sheaf of weight $dD_\mu$. This descent/subquotient formalism is BBD~5.1.2; purity of the summand is BBD~5.3.1. On its own smooth open Schubert stratum inside its support, $\IC_\lambda$ is the constant sheaf shifted by $D_\lambda$. Restriction to this open preserves both weight inequalities because $j^*=j^!$ (BBD~5.1.14). The Hom module shifted by $D_\lambda$ is therefore pure of weight $dD_\mu$, which means the module itself has weight $dD_\mu-D_\lambda$ by the shifted convention of BBD~5.1.8. This proves the all-eigenvalue purity statement without assuming arithmetic semisimplicity. The tensor equivalence identifies its dimension with the required multiplicity.

For the product group, external product gives the simple object corresponding to $V\boxtimes W$. Its weight is $(m-1)+(n-1)$. The same Schur-idempotent construction corresponds to $S_\nu(V\otimes W)$ as a $\GL_m\times\GL_n$ representation. The definition of the Kronecker coefficient and the just-proved Hom-weight calculation give the stated formula.

Finally consider $\GL_2$ and two minuscule steps. The twisted convolution variety is a smooth surface. Over the central lattice $t\mathcal O^2$, its fiber consists of the intermediate lattices between $t\mathcal O^2$ and $\mathcal O^2$, hence of lines in $\mathcal O^2/t\mathcal O^2$; it is $\PP^1$. The contribution to the perverse degree-zero stalk is $H^2(\PP^1)=\overline\Q_\ell(-1)$. To identify it as the determinant multiplicity, Satake gives
\[
 V\otimes V=\Sym^2V\oplus\wedge^2V,
\]
with determinant multiplicity one. Equivalently the non-open relevant fiber supplies this unique one-dimensional summand; the IC of the open two-dimensional stratum has no perverse degree-zero stalk contribution at the boundary point by the intermediate-extension condition. The corrected symmetric commutativity corresponds to the ordinary tensor swap, which sends
$e_1\otimes e_2-e_2\otimes e_1$ to its negative. The sole nonzero basis vector of a one-dimensional permutation representation would have to be fixed, so no such basis exists. The symmetric projector kills this line, while the alternating projector keeps it.
\end{proof}
This is a constructive reduction of the ordinary-geometry gate for \emph{classical} coefficients. Its purity uses the existing Weil/BBD theorems; it is not a new application of \fid{F03}, and it does not provide the path-hypersurface identities for a nonstandard canonical coefficient. The remaining problem is more specific: describe the actual projector image, its graded basis and its compatibility with the designated quantum operations, together with effective labels when a counting rule is required. The $\GL_2$ example prevents replacing that task by unsigned orbit-counting on the unprojected component basis.

The construction also provides an exact Grothendieck-framework dictionary. Convolution is a proper constructible pushforward; the Young idempotent is an actual projector in a characteristic-zero symmetric tensor category; its image is an effective object whose $K_0$ class is the Schur operation. The weights record the unnormalized IC shifts, so no half-Tate convention is hidden. The motivating language of projectors occurs in cote~16, scans~35, 37 and~39; the actual theorem here uses the modern Satake equivalence and the precisely cited SGA/BBD weight formalism, not an attribution to those leaves. The explicit symmetric-group action is the additional datum beyond virtual $\lambda$ identities.

\paragraph{\fid{F20}, motivic refinement: scalar Tate Frobenius.}
In the rational motivic Satake category over $\F_q$, these same Schur
projector images have coefficient motives
\[
 \mathbf1(-h)^{\oplus a},\qquad
 h=\frac{dD_\mu-D_\lambda}{2}
   =n(\lambda)-d\,n(\mu),\qquad
 n(\alpha)=\sum_i(i-1)\alpha_i,
\]
where $a=[S_\lambda V:S_\nu(S_\mu V)]$. For the product-group Kronecker
construction the coefficient is
\[
 \mathbf1\bigl(-n(\lambda)-n(\mu)\bigr)^{\oplus g(\lambda,\mu,\nu)}.
\]
Consequently the natural Frobenius in the \emph{specified} Weil
multiplicity spaces is the scalar $q^h$ and has no Jordan part.

\begin{proof}
The input is Richarz--Scholbach \cite[Theorem~4.1; Proposition~4.8;
Proposition~5.9; Definition~6.2; Lemma~6.3; Corollary~6.4;
Lemma~6.5]{RS}. These provide pure intersection motives, weight-exact
convolution, symmetric constraints compatible with the corrected
\(\ell\)-adic realization, and a semisimple abelian compact Satake
category generated by the $\IC_\lambda(n)$ and stable under convolution.
Thus the primitive rational Young idempotent acts motivically and its
image is pure of weight $w=dD_\mu$. Its simple summands have the form
$\IC_\lambda(n)$, of weight $D_\lambda-2n$. Purity forces
$n=(D_\lambda-w)/2=-h$. Compatibility with geometric Satake identifies
their multiplicity with $a$. Restriction of the corresponding isotypic
summand to its open stratum and removal of the shift $D_\lambda$ gives
exactly $\mathbf1(-h)^{\oplus a}$; Lemma~6.3 gives the coefficient-motive
Hom identification. Since $|\lambda|=d|\mu|$ and
$D_\alpha=(r-1)|\alpha|-2n(\alpha)$, the stated integer formula for
$h$ follows. The product-group proof replaces $D_\lambda$ by
$D_{\lambda,m}+D_{\mu,n}$ and $w$ by $d(m+n-2)$, giving
$h=n(\lambda)+n(\mu)$. Finally the realization of $\mathbf1(-h)$ is
$\overline\Q_\ell(-h)$, whose geometric Frobenius is exactly $q^h$.
\end{proof}
Here the coefficient motive is not the untwisted vector space
$\Hom_{\rm mot}(\IC_\lambda,P)$, which can vanish when $h\ne0$.
The relevant multiplicity vector space is
$\Hom_{\rm mot}(\IC_\lambda(-h),P)$; tensoring it with the Tate line
recovers the Weil multiplicity module. This is stronger than purity for
these natural classical objects, and uses the modified dual-group
weight coordinate of \cite[Theorem~6.8]{RS}. It neither chooses a basis
of the coefficient vector space nor computes its rank efficiently.
Arithmetic semisimplicity for arbitrary pure Weil modules is still
false, as in \fid{F05}.
